% !TeX root = ../queens.tex
\section{Reproducing the computation}\label{app:artifacts}
The programs described here are in the companion repository
\url{https://github.com/boonsuan/queens}. Its folder
\texttt{verification} contains the finite verification of
Sections~\ref{sec:local}--\ref{sec:finite}, the folder \texttt{oeis}
the checks for Section~\ref{sec:oeis-consequences}, and the folder
\path{fast_generator} the generator of Section~\ref{sec:generation}.
Each folder has a README with further details. The Python programs
need Python~3.10 or later and no other packages, and all
mathematical checks use integer arithmetic.

\subsection{The finite verification}\label{app:verification}
From the folder \texttt{verification}, run
\begin{verbatim}
python verify_tuples.py
python verify_bitmasks.py
python compare_verifiers.py
python check_correspondence.py
python test_rejection.py
\end{verbatim}
Each takes a few seconds.

The first command carries out the check of
Section~\ref{sec:closure-check}, including the starting checks of
Section~\ref{app:seed}. It uses the module \texttt{calculation.py},
which stores the eight fields of \eqref{eq:state} as a tuple and
implements Algorithm~\ref{alg:local-successors} with generators that
yield one result per branch. The second command
performs the same check independently: it shares no code with the
first, packs the records into integers, splits each symbol into its
two bits, and handles branching with explicit work lists. The third
converts both to a common encoding and checks that they reach the same
states with the same complete sets of successors. This agreement
checks the implementations against each other; the correspondence
with the actual infinite sequence is the induction in
Section~\ref{sec:identification}.

The fourth command computes $3000$ queens directly from
\eqref{eq:greedy}. For each column $30\le n<3000$, it runs the actual
branch of Lemma~\ref{lem:local-exactness}, and checks
that every requested symbol labels a history-graph edge and has index
at most $m+6<n$, and that the branch yields the actual next queen and
the actual records before column $n+1$. This is a finite consistency
test; the infinite conclusion rests on the verification and the
induction. The last command removes the edge carrying $\sigma_{30}$
from a copy of the history graph and confirms that both verifiers
reject it.

The command \texttt{python trace.py 41 44} prints the actual records
before each column in a range, the requests made, the queen chosen,
and the new symbol, checking each step against the board. It
reproduces the examples of Section~\ref{sec:worked-transition}.

The two verifiers report
\begin{center}
\begin{tabular}{lr}
\toprule
Quantity&Result\\
\midrule
History-graph vertices and edges&$2092$, $2603$\\
State-graph vertices and edges&$7014$, $8327$\\
Completed local choices&$7612$\\
Completed lower choices&$3153$\\
Largest fresh input offset&$5$\\
Largest row advance&$5$\\
Range of $w-r-|D|$ at lower choices&$[-4\dts2]$\\
Requests with no outgoing edge&$0$\\
\bottomrule
\end{tabular}
\end{center}
The state-graph counts refer to distinct states and ordered pairs of
states. The row search can branch again after a choice has been made,
so completed choices and state edges have different counts. The
observed maxima for fresh input and row advance are smaller than the
bounds established in Lemma~\ref{lem:local-exactness}. Only the
diagonal discrepancy bound $4$ is used in the main theorem.
Section~\ref{sec:sharper} also uses the upper end $2$
of the range of $w-r-|D|$.

The command \texttt{python sharper\_constants.py} repeats the
exploration of the first verifier and checks the facts used in
Section~\ref{sec:sharper}: every state satisfies
$-4\le w-|D|-\phi|R|\le1$, every lower choice has $w-r-|D|\le2$, and
the columns before $30$ satisfy $-3/\phi<\varepsilon(n)<2/\phi$ and
$d_j-j\le2$. The comparisons with $\phi$ are exact, made with integers
by squaring. Over the $7014$ states, $w-|D|-\phi|R|$ attains both
$-4$ and $1$.

The history graph is stored in \texttt{history.json} as
\texttt{\{"memory": 12, "vertices": [[h, mask], \ldots]\}}. Each
record is one vertex: \texttt{h} is its twelve-symbol word in base
four, with the newest symbol in the least significant digit, and bit
$s$ of \texttt{mask} is set when an edge labeled $s$ leaves it. The
completed graph is a \emph{certificate}: finite proof data whose
required properties the verifiers check. The verification and the
induction of Section~\ref{sec:identification} together give the
diagonal discrepancy bound; the construction below only explains where
the certificate comes from.

\subsection{Constructing the history graph}\label{app:construction}
From the folder \texttt{verification}, run
\begin{verbatim}
python construct_history_graph.py
python history_length_experiment.py
\end{verbatim}
The first command carries out the construction of
Section~\ref{sec:finite-statement}. It computes the first thirty
queens directly from \eqref{eq:greedy} and runs the calculation, adding
each missing output edge instead of failing, until a complete
exploration adds no edge. It does not read \texttt{history.json}; it
reports each exploration and confirms that the result is identical to
\texttt{history.json}. The construction is not part of the proof; the
verifiers check its result independently.

The second command repeats the construction with histories of each
length $L$ from $4$ through $14$, keeping the same start, calculation,
and Condition~\ref{cond:state}. Shorter lengths are not tested:
because $z\ge-4$, the upper-queen tests can involve the upper-column
bits of columns $m-4,\dots,m-1$ (compare Section~\ref{sec:generation}),
so the input history must contain them. Lengths $4$ through $11$ fail,
because the construction reaches a successor violating
Condition~\ref{cond:state}; for lengths $5$ through $11$, the violation
found has $w=5$. Lengths $12$, $13$, and $14$ succeed, and each
resulting graph passes the check of Section~\ref{sec:closure-check}.
Thus twelve is the smallest successful length in this setting; this
does not establish a minimum over other encodings or bounds.

\subsection{Additional checks for the OEIS consequences}\label{app:oeis-checks}

The computations for Section~\ref{sec:oeis-consequences} are not
needed for Theorem~\ref{thm:main}. They reuse the programs of the
folder \texttt{verification} with longer histories. From the root of
the repository, run
\begin{verbatim}
python oeis/run_all.py
\end{verbatim}
It takes under a minute, prints one line for each statement checked,
and writes the graphs, languages, and witnesses to
\path{oeis/results}.

\medskip\noindent
\textbf{The refined certificate.}
Replace both twelve-symbol histories by histories of length $40$,
keeping the records, Condition~\ref{cond:state}, and the calculation
of Section~\ref{app:output}. Start immediately before column $80$,
where direct construction gives $m=51$, $d=32$, $U(m-1)=31$, $w=-3$,
$z=-2$, and $R=D=A=\varnothing$; the input history, queue, and output
history cover indices $11$--$50$, $51$--$79$, and $40$--$79$.
Lemma~\ref{lem:local-exactness} remains valid: omitted upper columns
lie further in the past, requests still end by $m+6$, and
$n-m=U(m-1)+z\ge12-4=8$, so its hypothesis remains $U(m-1)\ge12$,
not $40$. The induction of Section~\ref{sec:identification}
therefore applies once the refined graph passes the check of
Section~\ref{sec:closure-check}. The construction of
Appendix~\ref{app:construction}, with length $40$ and start
column $80$, gives a graph with $16876$ vertices and $17499$ edges.
With this graph fixed, the calculation reaches $29267$ states and
$30000$ directed state edges, and every check passes.

\medskip\noindent
\textbf{Graphs for the gap and run sequences.}
The output symbol of a state-graph edge is the last symbol of the
successor's $H_{\mathrm{out}}$. Starting from a state just after an
upper output, follow lower outputs to the next upper output, and
replace this path by a single edge labeled by its length. The result
is a finite graph whose walks include the gap sequence $g$ after
column $80$. Following edges labeled $1$ and then one edge labeled
$k>1$, and labeling the result $k-1$, gives the graph for $r$ used in
Corollary~\ref{cor:oeis-runs}. For each $c$, its edges labeled $c$
form an acyclic subgraph, and reverse-topological dynamic programming
lists the lengths of maximal runs bounded by unequal labels.
Similarly, the labels along paths from one $3$ to the next give the
return words. The part of the gap graph avoiding $3$ is acyclic, so
this language is finite and is computed exactly. For each return word,
the union of the languages at its possible endpoints contains every
word that can follow it; it has a single element for exactly $63$
words. The shortest path between two edges labeled $4$ has length
$71$, and the paths of that length all have the same labels.

\medskip\noindent
\textbf{Occurrence witnesses.}
The graphs give only upper bounds: attainment and nonfaithfulness
come from the actual sequence. A bitboard implementation of the
defining rule computes $20000$ queens, which contain a witness for
every value in Corollary~\ref{cor:oeis-runs} and every run that
begins before the refined graph applies (through column $95$ for
A275887). A million-queen prefix, generated by a separate program
that agrees with the bitboard calculation on its first $200000$
queens, contains every return word by gap index $108015$, and two
different successors of each word that is not faithful by gap index
$573517$. The file \path{oeis/AUDIT.md} separates these proved
consequences from observations that the checked graphs do not settle.

\subsection{Running the fast generator}\label{app:stream-generation}

The folder \path{fast_generator} contains the C11 generator of
Section~\ref{sec:generation}, the program that builds its table, the
checks, the benchmark tools, and the recorded measurements. From that
folder, with a C11 compiler, \texttt{make}, and Python~3.10 or later,
run
\begin{verbatim}
make
./build/queens_fast --count 1000000
make verify
make test
\end{verbatim}
The second command generates the first million rows and prints a
summary with the last row and a checksum of all rows; add
\texttt{-{}-emit} to print every column and row. The interface in
\path{src/queens_fast.h} also supports single rows, filled arrays, and
pauses between calls. The command \texttt{make verify} rebuilds the
table and the finite cases from the history graph and checks that they
are identical to the files used, then replays every local step and
every four-input path through a separate implementation of the local
calculation. The command \texttt{make test} compares every coordinate
through the first million queens with a calculation using full
occupancy arrays, whose first $3000$ rows are checked against the
greedy rule, and exercises the interface. These are finite tests; the
correctness of every row is
Proposition~\ref{prop:stream-generation}.

The runs behind Table~\ref{tab:c-generation} are recorded in
\path{results/ec2-native.json}, with the machine and build details in
\path{results/ec2-environment.json}. The command
\begin{verbatim}
python3 tests/check_ec2_results.py \
  --environment results/ec2-environment.json
\end{verbatim}
rechecks the recorded runs and recomputes the table, and
\path{bench/run_benchmark.py} runs a new campaign on a Linux machine.
The comparator \path{knuth/knuth_packed.c} is derived from Knuth's
\texttt{infty-queens} by \path{knuth/derive_knuth_packed.py}.

The check of Knuth's ranges at the end of Section~\ref{sec:contraction}
uses \path{src/scan_bounds.c}. With GCC or Clang, run
\begin{verbatim}
make build/scan_bounds build/knuth_packed
python3 tests/check_bounds.py
python3 bench/run_bounds_scan.py --count 100000000000
\end{verbatim}
The scanner checks both intervals and their union, and records the
extreme deviations and the queens attaining them. Comparisons too close
to decide by a rigorously bounded floating-point approximation are
resolved with exact integer arithmetic. The recorded result is
\path{results/knuth-bounds-1e11.json}, and \path{REPORT.md} describes
both measurements in full.
